\section{Compact-Field Supervision of 3DGS}
\label{sec:compact-supervision}

The optimizer never sees an image; it compares two point-evaluable functions.  The teacher is
the frozen field $\field_v$ of \cref{eq:field}.  The student is the current 3D Gaussian set,
rendered at arbitrary image points by a sparse \emph{point rasterizer}.

\subsection{Point-sampled student rendering}
\label{sec:point-rasterizer}

Given sample locations $\bm{x} \in \Omega_v$, the point rasterizer projects the 3D Gaussians
into view $v$ exactly as the dense renderer does (EWA projection
$P_{iv} = J_{iv} \Sigma_i J_{iv}\T + 0.3\,I_2$, including the $0.3$-pixel dilation), sorts by
view depth, and alpha-composites front to back at each requested point:
\begin{equation}
  \student_v(\bm{x}) \;=\; \sum_{i} T_i(\bm{x})\, \alpha_i(\bm{x})\, \bm{c}_i(\bm{d}_{iv}),
  \qquad
  T_i(\bm{x}) = \prod_{i' < i} \big(1 - \alpha_{i'}(\bm{x})\big),
  \label{eq:student}
\end{equation}
with $\alpha_i(\bm{x}) = o_i \exp(-\tfrac12 q_{iv}(\bm{x}))$, $q_{iv}$ the Mahalanobis form of
$P_{iv}$, and view-dependent SH color $\bm{c}_i(\bm{d}_{iv})$.  The point path matches the
dense CPU reference renderer on color, alpha, depth, global visible order, five 3D
parameter-gradient families, and retained screen-space gradients within sealed tolerances, on
both synthetic grids and a calibrated 4{,}096-point replacement draw; this parity is the
correctness anchor that lets us reason about sampled training as subsampled dense
training.\evidence{ara/logic/claims.md C18;
benchmarks/results/20260716_point_rasterizer_parity_RESULT.json}
\toshow{Boundary: full-image calibrated parity, CUDA/gsplat point parity, and active off-grid
coordinate gradients are not yet established and must be either verified or explicitly scoped
out in the final text.}

\subsection{An unbiased sampled objective}
\label{sec:objective}

Let $p$ be the target risk measure on $\Omega_v$ (uniform over pixels, or a continuous area
measure), and let $q$ be a proposal distribution.  Training draws a \emph{fixed} number $A$ of
attempts $\bm{x}_1, \dots, \bm{x}_A \sim q$ per step and minimizes the importance-corrected
squared error with rejected attempts retained in the denominator:
\begin{equation}
  \widehat{L}_v \;=\;
  \frac{1}{A} \sum_{a=1}^{A}
  \mathbf{1}[\bm{x}_a \text{ accepted}]\;
  \frac{p(\bm{x}_a)}{q(\bm{x}_a)}\;
  \norm{\student_v(\bm{x}_a) - \field_v(\bm{x}_a)}_2^2 .
  \label{eq:sampled-loss}
\end{equation}
Keeping the denominator at the fixed attempt count $A$ (rather than the accepted count) is
what makes \cref{eq:sampled-loss} an unbiased estimator of the target risk under
rejection-mixture proposals; we verified the estimator identity by exact enumeration on finite
pixel grids, including invariance to attempt micro-chunking within
$2\times10^{-12}$.\evidence{ara/logic/claims.md C19}

Two evaluation risks are computed \emph{exactly} and are never replaced by the training loss:
$\Jpix$, the mean over all pixel centers, and $\Jarea$, a fixed four-offset quadrature per
pixel.  In experiment tables we report the validation risks $\Ju$ (uniform sampling) and
$\Jq$ (proposal-corrected), both against the frozen fields of held-out-for-validation
views.\evidence{benchmarks/compact_only_carrier_stage_ablation.py}

\paragraph{Proposal choice is settled empirically.}
Teacher-shaped Gaussian proposals are the obvious ``importance'' candidate; in a sealed
three-seed fixed-topology comparison they produced no material win --- the discrete-pixel
Gaussian mixture lost direction in all three seeds (geometric-mean AUC ratio $1.0246$), and
the continuous-area Gaussian mixture was directionally favorable in all seeds
($0.9911$) but missed the preregistered $0.95$ materiality floor.  The pipeline therefore
defaults to uniform proposals; this negative result is part of the method's evidence
base.\evidence{ara/logic/claims.md C20}

\subsection{Streaming supervision and the memory shape}
\label{sec:streaming}

Because both sides of \cref{eq:sampled-loss} are point queries, a training step touches:
the compact fields of the sampled views (kilobytes each), the sampled coordinate batch, the
projected student parameters, and the model/optimizer state.  No dense render target, no
image cache, and no all-view RGB tensor exist anywhere in the loop --- the fields are queried
on demand and the hard boundary of \cref{sec:arms} forbids materializing dense frames.  This
is the mechanism behind \cref{claim:vram}; what it buys is quantified (only) by the protocol
of \cref{sec:vram-protocol}.

\begin{ToShowBlock}[Supervision-side obligations]
(i) Freeze the Arm-1 optimizer configuration (sample budget per step, view schedule,
learning rates, topology policy, stopping rule) in the preregistration --- these are draft
blockers, deliberately not chosen in this text.
(ii) Show a sampled-versus-dense equivalence curve: quality of compact-supervised training as
a function of samples per step, against dense-RGB training at matched update counts
(\cref{fig:sample-budget}).
(iii) Establish or scope out CUDA point-rasterizer parity (the CPU reference anchors
correctness; throughput claims need the CUDA path).
\end{ToShowBlock}

\begin{figure}[t]
  \centering
  \figplaceholder{5cm}{%
    \textbf{Figure 4 --- sample-budget sweep.}  Held-out PSNR (y) versus samples per step
    (x, log scale) for compact-field supervision, with the dense-RGB control as a horizontal
    band and matched update counts.  One panel per scene; seeds as thin lines, median bold.
    Purpose: locate the budget where sampled supervision saturates, justifying the chosen
    default.  Data: Arm-1 training sweeps once the Arm-1 schedule is frozen.}
  \caption{\redcaption{Quality versus sample budget for compact supervision.  To be
    produced.}}
  \label{fig:sample-budget}
\end{figure}
